% Source: https://artofproblemsolving.com/community/c3h1383241_how_to_use_crt
% Stable identity: aops-topic-1383241-q01
\textbf{(a)} Modulo $b$: $c\,bb'\equiv0$ and $dd'\equiv1$, so $x_0\equiv a$. Modulo $d$: $a\,dd'\equiv0$ and $bb'\equiv1$, so $x_0\equiv c$.

If $x$ satisfies both congruences, then $b\mid x-x_0$ and $d\mid x-x_0$. Since $\gcd(b,d)=1$, this gives $bd\mid x-x_0$. Conversely, every $x\equiv x_0\pmod{bd}$ satisfies both congruences. $\blacksquare$

\textbf{(b)} With $b=5$ and $d=7$: $7\cdot3=21\equiv1\pmod5$, so $d'=3$; and $5\cdot3=15\equiv1\pmod7$, so $b'=3$. Then
\[x_0=2\cdot7\cdot3+3\cdot5\cdot3=42+45=87\equiv17\pmod{35}.\]
So $x\equiv\boxed{17}\pmod{35}$.

(For small moduli a quick alternative is to run through $3,10,17,\dots$, the numbers that are $3$ mod $7$, until one is $2$ mod $5$.)
